Multi-robot collision avoidance treats each robot as a point that can change velocity at any instant. A
trotting quadruped cannot: how well it absorbs a change of commanded velocity depends on where in the gait
cycle the change arrives. We measure this dependence for a B\'ezier-curve trot controller on a simulated
Unitree Go1 and use it in a decentralized swarm. The phase-aware layer solves reciprocal collision avoidance
(ORCA) one gait period ahead, releases each new velocity command at the gait phase our capability map
predicts is least destabilizing, and shapes it with a gait-aware slew limiter; urgent conflicts bypass the
timing. In MuJoCo, across 4 to 32 robots in antipodal swaps, pillar fields, a narrow gate and a 53\,m obstacle
course, the full system is compared with the same stack without timing (790 runs, paired by seed,
Holm-corrected). In crowded and obstructed scenarios it cuts urgent avoidance events by 22--39\% and
99th-percentile body tilt by 19--24\%, increases the closest robot--robot distance in formations by 2--5\,cm,
and had no falls in any run. Prediction without phase timing gives no consistent benefit, so the timing is the
active ingredient. In sparse traffic (four robots) the method costs more than it saves.
